% open-logic.tex
% برای تولید متن کامل با سبک توسعهٔ منطق باز کامپایل می‌شود

\documentclass[../include/open-logic-part]{subfiles}

\begin{document}

\clearpage

\begin{editorial}
این پرونده همهٔ محتوای گنجانده‌شده در پروژهٔ منطق باز را بارگذاری
می‌کند. اگر یادداشت‌های ویراستاری از این دست نمایش داده شوند، نشان
می‌دهند که این پرونده بی‌آنکه دربارهٔ چگونگی ارائهٔ این مطالب اندیشیده
شده باشد، کامپایل شده است. اگر این متن را می‌خوانید، احتمالاً
\emph{توصیه نمی‌شود} از این پروندهٔ پی‌دی‌اف برای تدریس یا مطالعه
استفاده کنید.

پروژهٔ منطق باز سازوکارهای بسیاری فراهم می‌کند که به کمک آنها می‌توان
متنی مناسب‌تر برای تدریس یا خودآموزی تولید کرد. برای نمونه، متن در
حالت پیش‌فرض همهٔ عملگرهای منطقی را اولیه می‌گیرد و در اثبات‌ها تمام
حالت‌های مربوط به همهٔ عملگرها را بررسی می‌کند. اما بسیار بهتر است
برخی از این حالت‌ها به‌صورت تمرین واگذار شوند. پروژهٔ منطق باز همچنان
در دست تکمیل است. برای ترغیب همکاری و بهبود، مطالب حتی هنگامی
گنجانده می‌شوند که صرفاً در حد پیش‌نویس‌اند، فاقد تمرین‌اند و مانند آن.
نسخهٔ پی‌دی‌افی که برای یک درس تهیه می‌شود، این بخش‌ها را در بر نخواهد
گرفت.

برای یافتن نسخه‌های پی‌دی‌اف مناسب‌تر برای تدریس و مطالعه،
\href{http://builds.openlogicproject.org/}{درس‌های نمونهٔ موجود در
  وبگاه پروژهٔ منطق باز} را ببینید. برای ساختن نسخهٔ خودتان، می‌توانید
کار را از
\href{https://github.com/OpenLogicProject/OpenLogic/tree/master/courses/sample}{پروندهٔ
  راه‌انداز نمونه} آغاز کنید یا برای دیدن نمونه‌های مفصل‌تر و
پیشرفته‌تر، کدهای مبدأ کتاب‌های درسیِ برگرفته از پروژه را بررسی کنید.
\end{editorial}

\olimport[sets-functions-relations]{sets-functions-relations-complete}

\olimport[propositional-logic]{propositional-logic}

\olimport[first-order-logic]{first-order-logic}

\olimport[model-theory]{model-theory}

\olimport[computability]{computability}

\olimport[turing-machines]{turing-machines}

\olimport[incompleteness]{incompleteness}

\olimport[second-order-logic]{second-order-logic}

\olimport[lambda-calculus]{lambda-calculus}

\olimport[many-valued-logic]{many-valued-logic}

\olimport[normal-modal-logic]{normal-modal-logic}

\olimport[applied-modal-logic]{applied-modal-logic}

\olimport[intuitionistic-logic]{intuitionistic-logic}

\olimport[counterfactuals]{counterfactuals}

\olimport[set-theory]{set-theory}

\olimport[methods]{methods}

\olimport[history]{history}

\olimport[reference]{reference}

\end{document}
